\subsection{正规坐标中的测地球体积}

Jacobi 场控制指数映射的体积伸缩；将这一伸缩积分，就得到曲率对小球体积的修正。

给定 $p\in M$，以测地距离 $r(x)=\operatorname{dist}_g(p,x)$ 定义球面看似天然。然而 $r$ 只在避开 $p$ 与割迹的区域光滑。指数映射 $\exp_p:T_pM\to M$ 在注入半径以内给出测地极坐标，度规具有形式
\[
 g=\dd r^2+h_{ab}(r,\theta)\dd\theta^a\dd\theta^b.
\]
局部坐标式不表示 $h_{ab}=r^2\gamma_{ab}$。后一个等式只在特别的几何中成立。

\begin{theorem}[测地球的面积和曲率修正]
\label{thm:ghs-small-ball-volume}
在 $p$ 的正规坐标 $x$ 附近，测度有展开
\begin{equation}\label{eq:ghs-normal-coordinate-measure}
 \sqrt{\det g(x)}=1-\frac16\operatorname{Ric}_{ij}(p)x^ix^j+O(|x|^3).
\end{equation}
因此半径 $r$ 的小测地球体积满足
\begin{equation}\label{eq:ghs-small-geodesic-ball}
 \operatorname{Vol}(B_r(p))=\omega_d r^d
 \left[1-\frac{\mathcal R(p)}{6(d+2)}r^2+O(r^4)\right],
\end{equation}
$\omega_d$ 是欧氏单位球体积；余项在足够光滑的固定度规下理解。
\end{theorem}
\begin{proof}
正规坐标中测地线为直线且 $\partial_k g_{ij}(p)=0$。四指标始终按前文定义排列：
\[
 R_{abcd}:=R(e_a,e_b,e_c,e_d)=g(R(e_a,e_b)e_c,e_d).
\]
它与先降低 $R^a{}_{bcd}$ 首指标的排列不同：
\[
 g_{ae}R^e{}_{bcd}=R_{cdba}=-R_{abcd}.
\]
由 Jacobi 场方程得到的正规坐标展开在这一排列下为
\[
 g_{ij}(x)=\delta_{ij}+\frac13R_{ikjl}(p)x^kx^l+O(|x|^3),
 \qquad \delta^{ij}R_{ikjl}=-\operatorname{Ric}_{kl}.
\]
例如常截面曲率 $k$ 时，$R_{ikjl}=k(\delta_{kj}\delta_{il}-\delta_{ij}\delta_{kl})$，故
\[
 g_{ij}(x)=\delta_{ij}-\frac{k}{3}(|x|^2\delta_{ij}-x_ix_j)+O(|x|^3).
\]
球面的 $k>0$ 给出负的切向面积修正，与测地线聚焦相符。
对矩阵 $I+A$ 使用 $\det(I+A)=1+\operatorname{tr}A+O(A^2)$ 以及平方根展开，得到 \cref{eq:ghs-normal-coordinate-measure}。
将其在半径 $r$ 的欧氏球上积分，利用各向同性
\[
 \int_{|x|<r}x^ix^j\dd^d x
 =\delta^{ij}\frac{\omega_dr^{d+2}}{d+2},
\]
二次修正变为 $-\mathcal R\,\omega_d r^{d+2}/[6(d+2)]$。常曲率时 $\mathcal R=d(d-1)k$，于是球面的相对体积修正为 $-d(d-1)kr^2/[6(d+2)]<0$。三阶角向奇函数积分为零，因此积分余项可提升到 $O(r^4)$ 的相对阶；这一步要求度规有充分阶数的光滑展开。
\end{proof}
